\documentclass[10pt]{article}
\usepackage[margin=0.72in]{geometry}
\usepackage{amsmath,amssymb,booktabs,array,tabularx,longtable,microtype,xcolor,hyperref,enumitem,fancyhdr}
\usepackage[T1]{fontenc}
\usepackage{lmodern}
\definecolor{navy}{HTML}{15304A}
\definecolor{teal}{HTML}{167D86}
\definecolor{pale}{HTML}{EEF5F6}
\definecolor{warn}{HTML}{8A4B08}
\hypersetup{colorlinks=true,linkcolor=navy,urlcolor=teal,citecolor=navy}
\setlist{nosep,leftmargin=1.3em}
\setlength{\parskip}{4pt}
\setlength{\parindent}{0pt}
\pagestyle{fancy}
\fancyhf{}
\lhead{\small Multiscale PC Cap for ePC Transformers}
\rhead{\small Design hypothesis --- 30 July 2026}
\cfoot{\thepage}
\newcommand{\obs}{\textbf{Published fact.}}
\newcommand{\hyp}{\textbf{Hypothesis.}}
\newcommand{\guess}{\textbf{Engineering guess.}}
\newcommand{\gate}[1]{\texttt{#1}}
\newcommand{\R}{\mathbb{R}}

\title{\vspace{-1.5em}\textbf{\color{navy}A Multiscale, Laterally Connected Predictive-Coding Cap\\for Continual Learning on ePC Transformers}\\[0.45em]
\large Architecture hypothesis, comparison with backpropagation, and coding-agent implementation guide}
\author{ZeroBot / ProtoCosmoBot\\\small Synthesis of the Bot Philosophy discussion with Ben Goertzel and ProtoMegaBot}
\date{30 July 2026}

\begin{document}
\maketitle
\begin{abstract}
This note proposes a small recurrent predictive-coding (PC) ``cortical island'' above an
error-predictive-coding (ePC) transformer.  Three cap levels combine hierarchical prediction,
bottom-up error, sparse lateral attractor dynamics, multiple learning timescales, and
recurrence-gated consolidation.  Novel input is first quarantined in a fast episodic level;
only repeated, coherent novelty is promoted into slower contextual and schema levels and
allowed to modulate base-network plasticity.  The main claim is conditional and falsifiable:
if an ePC and a backpropagation (BP) transformer have matched base-task quality, an ePC base
should provide a cleaner substrate for this cap because prediction errors, precision,
top-down replay, settling dynamics, and local plasticity modulation share one computational
idiom.  The report gives interfaces, update equations, failure controls, ablations, and
explicitly speculative starting configurations at GPT-2, GPT-3, GLM-5.2, and Kimi-K3 scale.
\end{abstract}

\colorbox{pale}{\parbox{0.96\linewidth}{
\textbf{Epistemic status.} This is a design hypothesis, not an established continual-learning
result.  Published model dimensions are labeled as facts; cap sizes and learning rates are
engineering guesses to seed sweeps.  In particular, ``one unified energy'' requires compatible
objectives and reciprocal parameterizations; merely joining two modules called PC does not
guarantee a global Lyapunov function or convergence.
}}

\section{Operational objective}
Let a base transformer produce a residual representation $z_t\in\R^d$ for a token, segment,
or pooled event at time $t$.  We want a cap that learns online while reducing destructive
interference, without requiring long globally synchronized retraining cycles.  It should:
\begin{enumerate}
  \item recognize familiar instances, novel combinations, and genuinely novel regimes;
  \item form sparse, reusable context attractors from repeated experience;
  \item preserve uncertain novelty in a fast buffer before altering consolidated knowledge;
  \item send the base a context vector, uncertainty, and spatially selective plasticity gates;
  \item support latent generative replay of old contexts;
  \item learn with bounded local updates and short, optionally asynchronous settling.
\end{enumerate}

The cap is not expected to solve continual learning alone.  Its value must be tested against
matched recurrent memory, replay, regularization, and parameter-isolation controls.

\section{Architecture: a three-level recurrent PC cap}
\subsection{State and connectivity}
Use three cap states $y_t^1,y_t^2,y_t^3$:
\[
z_t \leftrightarrow y_t^1 \leftrightarrow y_t^2 \leftrightarrow y_t^3 ,
\qquad
y_t^\ell \leftrightarrow L_\ell y_t^\ell ,
\qquad
y_{t-1}^\ell \rightarrow y_t^\ell .
\]
The arrows between levels are reciprocal but need not be transposes.  Each level contains:
\begin{itemize}
  \item a top-down generator $g_\ell$ and bottom-up error encoder $q_\ell$;
  \item fixed or slowly learned lateral inhibition for sparse competition;
  \item learned block-sparse lateral excitation for associative attractors;
  \item a recurrent trace for temporal persistence;
  \item a precision head $\pi^\ell$ and consolidation statistic $c^\ell$.
\end{itemize}

\begin{center}
\begin{tabularx}{0.97\linewidth}{>{\bfseries}p{0.13\linewidth} X X X}
\toprule
Level & Primary content & Plasticity & Characteristic role\\
\midrule
$y^1$ episodic & local motifs and recent residual patterns & fast & capture/quarantine novelty; sparse instance code\\
$y^2$ contextual & recurring coalitions and situations & medium & context attractors; identify novel combinations\\
$y^3$ schema & regimes, tasks, domains & slow & consolidated prior; replay seed; base-gating policy\\
\bottomrule
\end{tabularx}
\end{center}

\subsection{Local energy and settling}
One implementable energy for a microbatch/window is
\[
\mathcal E =
\sum_{\ell=0}^{2}\frac12
\left\|\Pi_\ell^{1/2}
\left(y^\ell-g_{\ell+1}(y^{\ell+1})\right)\right\|^2
+\sum_{\ell=1}^{3}\lambda_\ell\Phi_\ell(y^\ell;L_\ell)
+\sum_{\ell=1}^{3}\rho_\ell\|y^\ell-A_\ell y_{t-1}^\ell\|^2 ,
\]
where $y^0=z$, $\Pi_\ell=\mathrm{diag}(\pi^\ell)$ is bounded precision,
$\Phi_\ell$ combines sparse competition and lateral consistency, and $A_\ell$ is a
temporal trace.  During cap inference, freeze weights and perform $K_\ell$ damped local
state steps:
\[
y^\ell \leftarrow y^\ell-\alpha_\ell
\operatorname{clip}\!\left(\nabla_{y^\ell}\mathcal E,c_y\right).
\]
Do not assume monotonic energy descent: test it.  Backtrack or reduce $\alpha_\ell$ when
$\mathcal E_{k+1}>\mathcal E_k+\epsilon_E$.  For streaming use, layers may update
asynchronously using versioned state snapshots; stale reads must be bounded and measured.

\subsection{Lateral connections: two distinct functions}
The phrase ``lateral connection'' should not conceal two mechanisms:
\begin{enumerate}
  \item \textbf{competition/normalization:} fixed inhibitory or divisive dynamics produce
  sparse codes and prevent every attractor from co-activating;
  \item \textbf{association/precision:} learned block-sparse excitation represents
  co-occurring coalitions and contributes to confidence/precision.
\end{enumerate}
Precision must be bounded, e.g.
\[
\pi^\ell=\pi_{\min}+(\pi_{\max}-\pi_{\min})
\sigma(P_\ell y^\ell+b_\ell),
\]
because unconstrained learned precision can ``explain away'' errors by driving gain to zero
or infinity.  Precision is not identical to attractor depth; the proposed relationship is a
hypothesis requiring calibration against held-out prediction reliability.

\section{Continual-learning rule}
\subsection{Do not equate surprise with permission to overwrite}
Raw high error may be coherent novelty, an outlier, corruption, or a base-representation
shift.  Therefore use a two-stage gate:
\begin{enumerate}
  \item \textbf{capture:} large lower-level error creates or updates a fast $y^1$ trace while
  base plasticity is conservative;
  \item \textbf{promotion:} recurrence across independent windows, predictive improvement,
  and agreement at $y^2/y^3$ increase consolidation confidence; only then may the cap
  selectively open base plasticity.
\end{enumerate}
For event $u$, a simple promotion score is
\[
s_u=\sigma\!\left(a\log(1+n_u)+b\,\Delta\mathcal E_u
-c\,\mathrm{Var}(e_u)-d\,o_u-\tau\right),
\]
where $n_u$ is recurrence count, $\Delta\mathcal E_u$ is predictive improvement, and $o_u$
is an outlier score.  Require recurrence in separated windows, not repeated tokens within one
sequence.

\subsection{Multi-level error pattern and action}
\begin{center}
\begin{tabularx}{0.98\linewidth}{p{0.25\linewidth} p{0.28\linewidth} X}
\toprule
Observed cap error & Interpretation (tentative) & Plasticity response\\
\midrule
$e^1$ low & familiar instance & near-zero base update; optional slow rehearsal\\
$e^1$ high; $e^{2,3}$ low & new instance of known context/schema & moderate update near readout or selected upper base blocks\\
$e^1$ low; $e^{2,3}$ high & familiar features in unfamiliar composition & update attention/composition modules; protect feature modules\\
all levels high, low outlier score, recurrent & coherent new regime & recruit cap capacity, then progressively open selected base modules\\
all levels high, nonrecurrent or high outlier score & noise/corruption/one-off & quarantine; no broad base update\\
\bottomrule
\end{tabularx}
\end{center}

The cap emits three outputs:
\[
c_t=D_c y_t^1,\qquad
u_t=D_u(e_t^1,e_t^2,e_t^3,\pi_t),\qquad
m_{b,t}\in[0,1]\quad(b=1,\ldots,B),
\]
where $c_t$ is a residual context, $u_t$ is calibrated uncertainty, and $m_{b,t}$ gates the
learning rule of base block $b$.  Initialize $D_c$ at zero so adding the cap is functionally
identity-preserving.

\subsection{Local weight updates and consolidation}
For reciprocal prediction matrix $W_\ell$, a PC-local candidate update is
\[
\Delta W_\ell = \eta_\ell\,s_t\,m_{\ell,t}\,
\operatorname{clip}\!\left((\Pi_{\ell-1}e^{\ell-1})(y^\ell)^\top,c_W\right)
-\eta_\ell\lambda_W W_\ell .
\]
Use eligibility traces if updates are delayed.  Keep three learning-rate bands with
$\eta_1:\eta_2:\eta_3$ initially around $100:10:1$.  This ratio is a sweep seed, not a law.
Consolidation may also maintain an exponential moving average (EMA) teacher for each cap
level.  Replay begins at sampled $y^3$ schemas, settles laterally at each level, generates
pseudo-$z$, and interleaves it with new windows.  Reject replay samples whose energy,
precision, or diversity lies outside the training envelope.

\section{Why an ePC transformer is the cleaner base}
\hyp Conditional on comparable base-task quality and a correctly implemented ePC learner,
the cap should integrate more naturally with ePC than with a conventional BP transformer.
The reasons are structural, but none substitutes for a controlled comparison.

\begin{enumerate}
\item \textbf{Signal compatibility.}  Both modules expose layerwise predictions, errors,
precision, and reciprocal influence.  The interface can be one more PC boundary rather than
an auxiliary density model glued to a feed-forward representation.

\item \textbf{Surgical plasticity modulation.}  An ePC local update can directly multiply
its synapse/block error term by $m_{b,t}$.  BP can also be masked or modulated, so ``BP cannot
do this'' would be false; the narrower claim is that BP's downstream credit assignment and
optimizer state make the intervention more globally entangled and less PC-native.

\item \textbf{Native top-down influence and latent replay.}  A reciprocal ePC base already
contains a path through which cap hypotheses can constrain lower activities and reconstruct
old latent regimes.  A standard causal transformer needs adapters, prefix/context injection,
an auxiliary decoder, or input-space replay.

\item \textbf{Settling trajectory as evidence.}  Residual error, convergence time, energy
slope, and oscillation supply familiarity/uncertainty features beyond a single forward state.
These signals are useful only if calibrated; slow settling can also simply indicate poor
conditioning.

\item \textbf{Potential common objective.}  If the cap and ePC base share explicitly
compatible local energies and reciprocal mappings, joint inference may minimize a declared
global energy.  This enables stability analysis.  It is not automatic: asymmetric feedback,
stale asynchronous states, precision learning, and separate task losses can break the claim.

\item \textbf{One hierarchy with a timescale gradient.}  The cap/base boundary can become
an engineering seam within a deeper PC hierarchy: fast task-sensitive lower levels and slow
context/schema upper levels.  Continual learning may then emerge from sparse recruitment,
precision, replay, and timescale separation rather than a single anti-forgetting penalty.
\end{enumerate}

\subsection{Honest BP counterargument}
BP transformers have mature kernels, pretrained weights, scaling evidence, and effective
continual-learning tools: replay, adapters, routing, regularization, gradient projection, and
parameter isolation.  BP permits per-block gradient masks, so the comparison must not use a
straw man.  ePC additionally pays settling latency and may suffer equilibrium, stability, and
credit-approximation errors.  The defensible claim is therefore:
\[
\textit{At matched base quality, parameter budget, replay budget, and wall-clock, does the
PC cap exploit an ePC base more effectively?}
\]

\section{Implementation contract for coding agents}
\subsection{Module boundaries}
\begin{verbatim}
BaseProbe:
  observe(hidden_by_block, error_by_block, settle_trace) -> BaseObservation
CapStateStore:
  load(stream_id, version) / compare_and_swap(stream_id, state)
PCCap:
  settle(observation, prior_state, K, masks) -> CapInference
PromotionGate:
  score(cap_inference, recurrence_stats, outlier_stats) -> Promotion
PlasticityRouter:
  route(promotion, cap_errors, base_errors) -> block_gates
LocalUpdater:
  accumulate_eligibility(...) / apply(gates, learning_rates)
ReplaySampler:
  sample_schema() / decode_and_validate() -> pseudo_observation
Telemetry:
  record(errors, energy, precision, gates, promotions, interference)
\end{verbatim}
Keep these interfaces replaceable.  A first implementation may use ordinary autograd to
verify equations, but it must label that path \gate{reference\_autograd}; it must not be
reported as a local-learning implementation.

\subsection{Numerical invariants}
\begin{enumerate}
  \item Adding a zero-initialized cap changes base logits by at most declared tolerance.
  \item Base and cap weights remain frozen during state settling.
  \item Error, energy, and precision accumulation use FP32 even when activations are BF16.
  \item No dropout, mutable KV cache, or nondeterministic routing changes within a paired
  settling/replay comparison.
  \item Every state update is finite; precision stays in $[\pi_{\min},\pi_{\max}]$.
  \item Accepted settle steps satisfy the energy policy or record a backtrack/failure.
  \item Local updates have explicit stop-gradient boundaries and unit tests against analytic
  linear cases.
  \item Asynchronous state writes use version checks; rejected stale writes are counted.
  \item MoE routing metadata is recorded; cap gates must not silently change expert selection
  during a matched comparison unless routing modulation is the tested treatment.
\end{enumerate}

\subsection{Build sequence}
\begin{enumerate}
  \item Implement a linear three-level cap and verify analytic gradients, energy descent,
  identity initialization, and local-update direction.
  \item Add fixed lateral inhibition; test target sparsity and absence of dead units.
  \item Add learned block-sparse excitation with spectral-radius constraint.
  \item Add temporal traces and recurrence-based promotion; test outlier quarantine.
  \item Connect read-only to a frozen small ePC transformer; calibrate errors and precision.
  \item Enable cap-only learning, then gated adapters/LoRA, then selected base blocks.
  \item Add top-down replay only after reconstruction and diversity gates pass.
\end{enumerate}

\section{Quantitative starting hypotheses}
\subsection{What scale should track}
Cap width should track the base residual-stream width $d$ and diversity of online contexts,
not total model parameters.  This is crucial for sparse MoE systems: total expert parameters
do not all appear in one token's residual stream.  Start with cap widths
$c_1,c_2,c_3$, factorized reciprocal maps of ranks $r_1,r_2,r_3$, lateral blocks of 64 units,
and a small prototype bank.  Parameter estimates below include factorized reciprocal maps
plus one 64-wide lateral and one lighter temporal operator per level; they exclude optimizer
state, replay buffers, and optional decoders.

\begin{small}
\begin{longtable}{>{\raggedright\arraybackslash}p{0.13\linewidth}
>{\raggedright\arraybackslash}p{0.18\linewidth}
>{\raggedright\arraybackslash}p{0.19\linewidth}
>{\raggedright\arraybackslash}p{0.17\linewidth}
>{\raggedright\arraybackslash}p{0.19\linewidth}}
\toprule
Base case & Published base configuration used & Proposed cap widths $(c_1,c_2,c_3)$ &
Map ranks $(r_1,r_2,r_3)$ & Core cap estimate and starting regime\\
\midrule
\endfirsthead
\toprule
Base case & Published base configuration used & Proposed cap widths &
Map ranks & Core cap estimate and starting regime\\
\midrule
\endhead
GPT-2 small &
$d=768$, 12 layers, 12 heads, 124M class &
$(512,256,128)$ &
$(128,64,32)$ &
$\sim0.6$M parameters. $K=(3,3,5)$; 64--256 context prototypes; cap LR
$10^{-4},10^{-5},10^{-6}$; base initially frozen, then top 2--4 blocks via LoRA.\\
\addlinespace
GPT-3 175B &
$d=12{,}288$, 96 layers, 96 heads &
$(4096,2048,1024)$ &
$(512,256,128)$ &
$\sim22$M core. $K=(2,3,4)$; 1K--4K prototypes; cap LR
$3{\times}10^{-5},3{\times}10^{-6},3{\times}10^{-7}$; gate 8--12 block groups rather
than 96 independent scalars initially.\\
\addlinespace
GLM-5.2 &
$d=6144$, 78 layers; 256 routed experts, top-8; first 3 MLP layers dense &
$(3072,1536,768)$ &
$(384,192,96)$ &
$\sim10$M core. $K=(2,3,4)$; 1K--2K prototypes. Separate gates for attention,
shared expert, routed experts, and router; freeze router first.\\
\addlinespace
Kimi-K3 &
text $d=7168$, 93 layers; 896 experts, top-16 plus 2 shared; hybrid KDA/full-attention &
$(3584,1792,896)$ &
$(448,224,112)$ &
$\sim14$M core. $K=(2,3,4)$; 2K--4K prototypes. Gate KDA, full-attention,
shared experts, routed experts, and router separately; begin text-only and freeze router.\\
\bottomrule
\end{longtable}
\end{small}

\obs The GLM-5.2 values above come from the official public configuration:
hidden size 6144, 78 layers, 256 routed experts, 8 selected per token, one shared expert, and
three initial dense MLP layers.  \obs The Kimi-K3 values come from its official public
configuration/model card: text hidden size 7168, 93 layers, 896 experts, 16 selected per
token, two shared experts, and a hybrid of KDA and full-attention layers.  \guess All cap
dimensions, ranks, prototype counts, settling steps, and learning rates are guesses.

\subsection{Why not scale the cap linearly with total parameters?}
At GPT-3 and MoE scale, dense $d\times c_1$ reciprocal maps are unnecessarily large.
Factorize them and share cap computation across block groups.  The cap represents slow
contexts over the residual stream; it need not mirror all base layers or experts.  For MoE,
start with residual/context gating and router readout only.  Allow router plasticity only after
showing that context gating improves retention without routing collapse or expert starvation.

\subsection{Memory and cadence}
\begin{itemize}
  \item Run cap inference on pooled 8--64-token events rather than every token in the first
  prototype.  Compare mean, attention, and learned event pooling.
  \item Maintain fast traces for 32--256 events, prototypes for 64--4096 contexts depending
  on scale, and an explicit eviction/merge policy.
  \item Try one cap update per 8--32 base events; one $y^2$ consolidation per 8--32 $y^1$
  promotions; one $y^3$ update per 8--32 $y^2$ promotions.
  \item Limit replay initially to 10--25\% of online events and match wall-clock in controls.
\end{itemize}

\section{Tuning plan and diagnostics}
\subsection{Tune in this order}
\begin{enumerate}
  \item \textbf{State dynamics:} $\alpha_\ell$, damping, normalization, spectral radius,
  $K_\ell$, energy/backtracking rate.
  \item \textbf{Code quality:} sparsity target (start 5--15\%), dead-unit fraction, attractor
  recovery under corruption, context separation.
  \item \textbf{Precision calibration:} reliability diagrams and expected calibration error
  for predicted reconstruction/prediction quality.
  \item \textbf{Promotion:} recurrence threshold, separated-window requirement, outlier gate,
  false promotion/false quarantine.
  \item \textbf{Plasticity routing:} start cap-only, then adapters, then base block groups.
  \item \textbf{Replay:} only after pseudo-latent fidelity, diversity, and no-collapse tests.
\end{enumerate}

\subsection{Required telemetry}
Log per level: normalized error, energy before/after each settle step, precision quantiles,
state norm, sparsity, attractor assignment entropy, recurrence counts, promotions, merges,
evictions, and update norms.  Log per base block: gate, local error, update norm, representation
drift, router/expert statistics for MoE, and settling cost.  Report throughput, latency, memory,
and energy/backtrack failures.  Averages alone are insufficient; retain tails and per-context
breakdowns.

\subsection{Failure signatures}
\begin{center}
\begin{tabularx}{0.98\linewidth}{>{\bfseries}p{0.24\linewidth} X X}
\toprule
Signature & Likely cause & First response\\
\midrule
all contexts merge & weak inhibition, high lateral radius, overly fast $y^2/y^3$ learning &
increase competition; lower LR/radius; add orthogonality/diversity control\\
every event becomes new & miscalibrated precision, base drift, threshold too low &
normalize errors by context; slow adapter; raise separated recurrence threshold\\
no promotion & cap too rigid or precision too low & lower threshold; increase $y^1$ capacity;
inspect predictive-improvement term\\
energy oscillates & step too large, asymmetric maps, stale asynchronous state &
damp/backtrack; bound delay; test reciprocal parameterization\\
old tasks retained, new task not learned & gates too conservative & open adapters/upper blocks;
reduce consolidation delay\\
new task learned, old tasks collapse & outlier/surprise directly opens base & enforce quarantine;
increase replay or parameter isolation\\
MoE expert starvation & cap/router feedback collapse & freeze router; entropy/load-balance gate;
separate shared/routed-expert modulation\\
\bottomrule
\end{tabularx}
\end{center}

\section{Evaluation and falsification}
Use sequential regimes with controlled overlap: permutations, new compositions of old
features, coherent new domains, and injected outliers.  Report:
\[
\text{average accuracy/perplexity},\quad
\text{backward transfer},\quad
\text{forward transfer},\quad
\text{worst-case forgetting},\quad
\text{adaptation area-under-curve}.
\]
Also measure context identification, promotion precision/recall, error calibration, replay
coverage, base representation drift, and wall-clock/energy cost.

\textbf{Minimum matched arms:}
\begin{enumerate}
  \item ePC base alone;
  \item ePC + one-layer novelty cap;
  \item ePC + hierarchical cap without lateral excitation;
  \item ePC + full cap without replay;
  \item ePC + full cap with replay;
  \item BP + same-cap interface via adapters and masked gradients;
  \item BP + matched recurrent memory;
  \item BP + strong replay/regularization or parameter-isolation baseline.
\end{enumerate}
Match trainable parameters, replay bytes, online examples, optimizer state, and wall-clock
as separate comparisons.  A parameter-matched win that disappears under wall-clock matching
is not sufficient.

\textbf{Central falsifier.} Reject the strong architectural hypothesis if, across multiple
regime sequences and seeds, the ePC+cap fails to reduce forgetting/adaptation tradeoff relative
to the strongest matched BP baseline, or if any gain is explained by extra replay, parameters,
settling compute, or privileged task-boundary information.

\section{Recommended first experiment}
Start at GPT-2-small or smaller, not at GPT-3/MoE scale:
\begin{enumerate}
  \item frozen pretrained base; train only the cap to predict pooled top-layer residuals;
  \item three synthetic stream regimes plus outlier bursts;
  \item compare one-layer, three-layer no-lateral, and full block-lateral cap;
  \item pass gates for energy behavior, context recovery, precision calibration, and outlier
  quarantine;
  \item then permit zero-initialized context injection and LoRA on the top four blocks;
  \item finally compare ePC and BP bases under matched parameter and wall-clock budgets.
\end{enumerate}
Only after the small model establishes a stable mechanism should the cap be ported to large
dense or MoE bases.  At MoE scale, router modulation is a separate experiment, not a default.

\section{Sources and provenance}
\begin{itemize}
  \item Bot Philosophy Telegram discussion, messages 15198--15227, 30 July 2026:
  initial episodic cap, multilayer/lateral extension, and ePC-vs-BP rationale.
  \item OpenAI, \emph{Language Models are Unsupervised Multitask Learners} (2019), GPT-2
  model configurations: \url{https://cdn.openai.com/better-language-models/language_models_are_unsupervised_multitask_learners.pdf}
  \item Brown et al., \emph{Language Models are Few-Shot Learners} (2020), GPT-3 175B
  configuration: \url{https://arxiv.org/abs/2005.14165}
  \item Z.ai official GLM-5.2 configuration (retrieved 30 July 2026):
  \url{https://huggingface.co/zai-org/GLM-5.2/blob/main/config.json}
  \item Moonshot AI official Kimi-K3 configuration/model card (retrieved 30 July 2026):
  \url{https://huggingface.co/moonshotai/Kimi-K3/blob/main/config.json},
  \url{https://github.com/MoonshotAI/Kimi-K3}
\end{itemize}

\vfill
\colorbox{pale}{\parbox{0.96\linewidth}{
\textbf{Bottom line.} Build the cap as a small, slow, recurrent PC hierarchy that first
captures and tests novelty, then consolidates it.  The ePC advantage is not mystical:
it is the possibility of sharing error, precision, top-down inference, and locally gated
plasticity across one declared dynamical system.  Whether that coherence beats mature BP
continual-learning methods is an empirical question, and the implementation should be
organized to make a negative result as informative as a positive one.
}}
\end{document}
